%%%% LISTA DE ABREVIATURAS E SIGLAS
%%
%% Relação alfabética das siglas. A lista impressa sai de \incluirlistadeacronimos{glossaries}:
%% só entra sigla citada com \gls ou \glspl. A primeira citação no texto expande a forma longa;
%% a lista usa o campo description.
%% gRPC não é sigla neste trabalho e não tem entrada aqui.

%% Abreviaturas: \abreviatura{rótulo}{representação}{definição}
%\abreviatura{art.}{art.}{Artigo}

%% Siglas. O tipo "sigla" precisa estar na lista de acrônimos antes do primeiro \newacronym,
%% senão \gls não expande na primeira citação.
\DeclareAcronymList{sigla}

\newacronym[type=\pseudoacronymtype, sort={CPU},
  description={Unidade Central de Processamento, do inglês \textit{Central Processing Unit}}]
  {cpu}{CPU}{unidade central de processamento, do inglês \textit{Central Processing Unit}}

\newacronym[type=\pseudoacronymtype, sort={EBPF},
  description={Filtro de Pacotes Berkeley estendido, do inglês \textit{extended Berkeley Packet Filter}}]
  {ebpf}{eBPF}{filtro de pacotes Berkeley estendido, do inglês \textit{extended Berkeley Packet Filter}}

\newacronym[type=\pseudoacronymtype, sort={GRAPHQL},
  description={Linguagem de consulta e manipulação de dados, do inglês \textit{Graph Query Language}}]
  {graphql}{GraphQL}{linguagem de consulta e manipulação de dados, do inglês \textit{Graph Query Language}}

\newacronym[type=\pseudoacronymtype, sort={HPACK},
  description={Algoritmo de compressão de cabeçalhos do HTTP/2}]
  {hpack}{HPACK}{algoritmo de compressão de cabeçalhos do HTTP/2}

\newacronym[type=\pseudoacronymtype, sort={HTTP},
  description={Protocolo de Transferência de Hipertexto, do inglês \textit{Hypertext Transfer Protocol}}]
  {http}{HTTP}{protocolo de transferência de hipertexto, do inglês \textit{Hypertext Transfer Protocol}}

\newacronym[type=\pseudoacronymtype, sort={HTTP2},
  description={Versão 2 do protocolo de transferência de hipertexto, do inglês \textit{Hypertext Transfer Protocol}}]
  {http2}{HTTP/2}{versão 2 do protocolo de transferência de hipertexto, do inglês \textit{Hypertext Transfer Protocol}}

\newacronym[type=\pseudoacronymtype, sort={JWT},
  description={Token Web JSON, do inglês \textit{JSON Web Token}}]
  {jwt}{JWT}{token web JSON, do inglês \textit{JSON Web Token}}

\newacronym[type=\pseudoacronymtype, sort={L7},
  description={Camada 7 (aplicação) do modelo OSI}]
  {l7}{L7}{camada 7, de aplicação, do modelo OSI}

\newacronym[type=\pseudoacronymtype, sort={MTLS},
  description={TLS mútuo, em que cliente e servidor se autenticam por certificados, do inglês \textit{mutual Transport Layer Security}}]
  {mtls}{mTLS}{TLS mútuo, do inglês \textit{mutual Transport Layer Security}}

\newacronym[type=\pseudoacronymtype, sort={MTTR},
  description={Tempo médio de recuperação, do inglês \textit{Mean Time To Recovery}}]
  {mttr}{MTTR}{tempo médio de recuperação, do inglês \textit{Mean Time To Recovery}}

\newacronym[type=\pseudoacronymtype, sort={NOSQL},
  description={Não somente SQL, do inglês \textit{Not Only SQL}}]
  {nosql}{NoSQL}{não somente SQL, do inglês \textit{Not Only SQL}}

\newacronym[type=\pseudoacronymtype, sort={RBAC},
  description={Controle de acesso baseado em papéis, do inglês \textit{Role-Based Access Control}}]
  {rbac}{RBAC}{controle de acesso baseado em papéis, do inglês \textit{Role-Based Access Control}}

\newacronym[type=\pseudoacronymtype, sort={REST},
  description={Transferência de Estado Representacional, do inglês \textit{Representational State Transfer}}]
  {rest}{REST}{transferência de estado representacional, do inglês \textit{Representational State Transfer}}

\newacronym[type=\pseudoacronymtype, sort={RPC},
  description={Chamada de Procedimento Remoto, do inglês \textit{Remote Procedure Call}}]
  {rpc}{RPC}{chamada de procedimento remoto, do inglês \textit{Remote Procedure Call}}

\newacronym[type=\pseudoacronymtype, sort={SQL},
  description={Linguagem de Consulta Estruturada, do inglês \textit{Structured Query Language}}]
  {sql}{SQL}{linguagem de consulta estruturada, do inglês \textit{Structured Query Language}}

\newacronym[type=\pseudoacronymtype, sort={TCP},
  description={Protocolo de Controle de Transmissão, do inglês \textit{Transmission Control Protocol}}]
  {tcp}{TCP}{protocolo de controle de transmissão, do inglês \textit{Transmission Control Protocol}}

\newacronym[type=\pseudoacronymtype, sort={TLS},
  description={Segurança da Camada de Transporte, do inglês \textit{Transport Layer Security}}]
  {tls}{TLS}{segurança da camada de transporte, do inglês \textit{Transport Layer Security}}

\newacronym[type=\pseudoacronymtype, sort={TTL},
  description={Tempo de vida, do inglês \textit{Time To Live}}]
  {ttl}{TTL}{tempo de vida, do inglês \textit{Time To Live}}

\newacronym[type=\pseudoacronymtype, sort={VPS},
  description={Servidor Virtual Privado, do inglês \textit{Virtual Private Server}},
  longplural={servidores virtuais privados, do inglês \textit{Virtual Private Server}},
  shortplural={VPS}]
  {vps}{VPS}{servidor virtual privado, do inglês \textit{Virtual Private Server}}

%% Acrônimos: \acronimo{rótulo}{representação}{definição}
%\acronimo{gimp}{Gimp}{Programa de Manipulação de Imagem GNU, do inglês \textit{GNU Image Manipulation Program}}
